\documentclass[dvipdfmx,a4paper]{article}
\usepackage[margin=1in]{geometry}
\usepackage{amsmath,amssymb,amsthm,amsfonts}
\usepackage{booktabs}
\usepackage{array}
\usepackage{enumitem}
\usepackage[dvipdfmx]{hyperref}
\usepackage[nameinlink,noabbrev]{cleveref}
\usepackage{xcolor}
\usepackage{siunitx}

\newtheorem{theorem}{Theorem}[section]
\newtheorem{proposition}[theorem]{Proposition}
\theoremstyle{definition}
\newtheorem{definition}[theorem]{Definition}
\newtheorem{remark}[theorem]{Remark}

\title{Resonance Warp in PGA--DST:\\
Layer-Hopping vs On-Demand Phase Entry}
\author{https://github.com/hypernumbernet}
\date{\today}

\begin{document}

\maketitle

\begin{abstract}
We develop a hypothetical interstellar propulsion framework within the projective geometric algebra (PGA) extension of dual spacetime theory (DST). Gravity and inertia emerge from torsional mismatch between particle-intrinsic usual and dual spacetimes, lifted into $G(3,1,1)$ with a ten-dimensional bivector generator space (hyperbolic, cyclic, and null sectors). The extended motor $M=\exp(\Omega_{\mathrm{biv}}^{(5)})=RT$ is unitary ($M\widetilde{M}=1$), preserving causal propagation on the light cone $\eta_{\mu\nu}\lambda^\mu\lambda^\nu=0$. The torsion scalar $J$ changes sign at odd multiples of $\pi$ in the dual rotor angle $\theta_a$, yielding repulsive layers ($J<0$) that enable \emph{effective} superluminal coasting without local causality violation.

Unlike Alcubierre-type metric engineering, warp in PGA--DST is not about bending a spacetime fabric with exotic matter; it is about coherently tuning collective dual-rotor resonance and entering $J<0$ layers via circularly polarised THz driving. We quantitatively compare two operational regimes: (1)~\textit{layer-hopping}, exploiting slow cosmic phase drift of background dual rotors, and (2)~\textit{on-demand entry}, applying active phase driving. For million-tonne-class vessels at $0.1$--$0.5c$, layer-hopping reduces entry energy by 4--7 orders of magnitude, into the \si{kWh}--low \si{MWh} range. All derivations respect the $G(3,1,1)$ algebra, teleparallel equivalence with general relativity, and unitary motor causality. We connect repulsive $J<0$ layers, helicity-controlled gravity engineering, layer stability, and an experimental roadmap. This paper is explicitly \emph{hypothetical}: macroscopic dual-phase coherence remains unverified.
\end{abstract}

\section{Introduction}
\label{sec:intro}

Conventional warp-drive proposals, beginning with Alcubierre's metric bubble~\cite{alcubierre1994}, require exotic negative energy density and violate the null energy condition (NEC)~\cite{ford2000}. Engineering such configurations at macroscopic scales appears prohibitive under known physics.

Dual spacetime theory (DST) offers a fundamentally different ontology~\cite{dual2025}. There is no autonomous spacetime continuum; instead, \emph{space emerges as collective resonance} among dual rotors carried by individual particles. Gravity is torsional mismatch between the usual and dual sectors, reformulated as teleparallel gravity (TEGR)~\cite{tegr} with vanishing curvature and non-zero torsion. The PGA extension~\cite{pga2025} embeds this structure in projective geometric algebra $G(3,1,1)$, adjoining a null vector $e_4$ so that translations, light propagation, and torsion are unified in a single ten-dimensional bivector object propagated by a unitary motor.

In this framework, \textbf{warp propulsion} is not metric engineering but \emph{resonance tuning}: shifting the collective dual-rotor phase to enter repulsive torsional layers ($J<0$) where effective displacement exceeds the local light-speed bound, while information remains constrained to the light channel. This paper is a \emph{hypothetical} extension of PGA--DST to interstellar propulsion. We preserve the quantitative energy estimates of the earlier DST warp analysis while grounding them in the PGA algebraic structure, local generation of repulsive $J<0$ layers, and gravity-engineering mechanisms from the PGA textbook programme.

\section{PGA--DST Foundations}
\label{sec:foundations}

\subsection{Three-sector structure in G(3,1,1)}

Each massive particle carries paired compactified Minkowski structures encoded in $\mathrm{Cl}(3,1)\cong\mathbb{H}\otimes\mathbb{H}$. The PGA lift adjoins $e_4$ with $e_4^2=0$ and $\{e_4,e_\mu\}=0$, extending the algebra to $G(3,1,1)$ (32 real dimensions). Vector generators are
\[
e_0=j,\quad e_1=kI,\quad e_2=kJ,\quad e_3=kK,
\]
with signature $(-,+,+,+)$. The bivector space has dimension ten and decomposes into three sectors~\cite{pga2025}:

\begin{table}[ht]
\centering
\caption{Ten-dimensional generator structure in PGA--DST}
\label{tab:generators}
\begin{tabular}{@{}clcl@{}}
\toprule
Sector & Generators & Square & Role \\
\midrule
Hyperbolic & $iI,\, iJ,\, iK$ & $+1$ & Usual-spacetime boosts \\
Cyclic & $I,\, J,\, K$ & $-1$ & Dual-spacetime rotations \\
Null & $N_\mu=e_4\wedge e_\mu$ & $0$ & Translations \\
\bottomrule
\end{tabular}
\end{table}

The null sector realises the Poincar\'{e} translation ideal: $N_\mu N_\nu=0$ for all $\mu,\nu$, and the Lie algebra closes as $\mathfrak{iso}(3,1)=\mathfrak{so}(3,1)\oplus\mathfrak{so}(3,1)\ltimes\mathbb{R}^{3,1}$.

\subsection{Rotors, motors, and the duality map}

The usual and dual rotors are
\begin{align}
R_{\mathrm{usual}} &= \exp\!\Bigl(\sum_{a=1}^3 \frac{\phi_a}{2}\, i\Gamma_a\Bigr),
&
R_{\mathrm{dual}} &= \exp\!\Bigl(\sum_{a=1}^3 \frac{\theta_a}{2}\, \Gamma_a\Bigr),
\end{align}
where $\Gamma_a\in\{I,J,K\}$. The relative mismatch rotor is
\begin{equation}
\Omega = R_{\mathrm{usual}}^\dagger R_{\mathrm{dual}},
\qquad
\Omega_{\mathrm{biv}} = \log\Omega.
\label{eq:omega}
\end{equation}
The duality map $X\mapsto Xi$ (with $i=e_0e_1e_2e_3$) interchanges hyperbolic and cyclic sectors while leaving the null sector closed, since $e_4$ commutes with $i$.

The extended mismatch bivector and motor are
\begin{equation}
\Omega_{\mathrm{biv}}^{(5)}
= \underbrace{\sum_{a=1}^3 \Bigl(\tfrac{\alpha_a}{2}B_a^{+} + \tfrac{\beta_a}{2}B_a^{-}\Bigr)}_{\Omega_{\mathrm{torsion}}}
+ \underbrace{\sum_{\mu=0}^{3}\tfrac{\lambda^\mu}{2}N_\mu}_{\Omega_{\mathrm{trans}}},
\label{eq:omega5}
\end{equation}
\begin{equation}
M = \exp(\Omega_{\mathrm{biv}}^{(5)}) = R\,T,
\qquad
T = 1 + \sum_{\mu=0}^{3}\tfrac{\lambda^\mu}{2}N_\mu,
\label{eq:motor}
\end{equation}
where the null exponential truncates at first order because $N_\mu N_\nu=0$.

\begin{theorem}[Motor unitarity]
\label{thm:unitarity}
Every bivector generator is reversal-odd ($\widetilde{B}=-B$). Consequently $R\widetilde{R}=1$, $T\widetilde{T}=1$, and $M\widetilde{M}=1$. The sandwich product $X'=MX\widetilde{M}$ preserves the degenerate $G(3,1,1)$ metric and causal structure.
\end{theorem}

\subsection{Torsion scalar and gravitational action}

The torsion scalar is defined via the Cartan--Killing form on the six-dimensional torsion sector:
\begin{equation}
J = \tfrac{1}{16}\,B_{\mathrm{Killing}}(\Omega_{\mathrm{biv}},\Omega_{\mathrm{biv}})
= \tfrac{1}{2}\sum_{a=1}^{3}(\alpha_a^2 - \beta_a^2).
\label{eq:J}
\end{equation}
The hyperbolic generators contribute $B(B_a^{+})=+8$ and the cyclic generators $B(B_a^{-})=-8$, so $J$ changes sign when dual-sector rotation dominates. In the small-angle regime, with $\alpha_a\sim\sinh(\phi_a/2)$ and $\beta_a\sim\sin(\theta_a/2)$, \eqref{eq:J} reduces to the finite-angle form
\begin{equation}
J \approx \sum_{a=1}^{3}\Bigl[\sinh^2\!\Bigl(\tfrac{\phi_a}{2}\Bigr) - \sin^2\!\Bigl(\tfrac{\theta_a}{2}\Bigr)\Bigr],
\label{eq:J-finite}
\end{equation}
used in the operational estimates below.

The PGA extended invariant couples torsion to translation:
\begin{equation}
J^{(5)} = J + \tfrac{1}{2}\eta_{\mu\nu}\lambda^\mu\lambda^\nu,
\label{eq:J5}
\end{equation}
where $\eta_{\mu\nu}=\mathrm{diag}(-1,+1,+1,+1)$. The translational term vanishes on the light cone $\eta_{\mu\nu}\lambda^\mu\lambda^\nu=0$, embedding null propagation algebraically.

Gravity follows from the action
\begin{equation}
S = \frac{c^4}{16\pi G}\int J\,\sqrt{-g}\,d^4x,
\label{eq:action}
\end{equation}
which is equivalent to the Einstein--Hilbert and TEGR actions: all GR solutions, including Schwarzschild, are recovered from tetrads generated by unitary motors, without invoking curvature~\cite{pga2025,tegr}. Interpretation changes; equations of motion do not.

\section{What Warp Means in PGA--DST}
\label{sec:warp-ontology}

\subsection{Space as collective resonance}

In PGA--DST, macroscopic space is not a pre-existing manifold but the statistical alignment of particle-intrinsic rotors across $\sim 10^{32}$ baryons per tonne of matter. The scalar $J$ measures misalignment. When neighbouring dual rotors are aligned ($\theta_a\approx 2n\pi$), $J>0$ and we experience attractive gravity. When driven to repulsive points ($\theta_a\approx(2n+1)\pi$), collective resonance shifts and effective geometry changes.

Warp propulsion takes two distinct operational meanings:

\begin{enumerate}[leftmargin=2em]
  \item \textbf{Resonance-shifting warp (layer-hopping).}
  The vessel remains within the collective resonance manifold but shifts local dual-rotor phase to enter a repulsive layer ($J<0$). The spacecraft rides a different harmonic of the same resonant structure. This is the primary mechanism analysed below.

  \item \textbf{Subspace-transition warp (dual-sector immersion).}
  The vessel actively decouples from ambient collective resonance and immerses primarily in the dual sector. Re-emergence requires precise re-synchronisation with the target region's phase. This is energetically far costlier than layer-hopping because it breaks macroscopic coherence entirely.
\end{enumerate}

Thus warp is \emph{changing the tuning of cosmic resonance}, not bending or tearing a fabric. The engineering challenge shifts from exotic negative energy to coherent dual-phase control at macroscopic scales.

\subsection{Comparison with Alcubierre warp drives}

\begin{table}[ht]
\centering
\caption{Alcubierre metric engineering vs.\ PGA--DST resonance warp (hypothetical)}
\label{tab:alcubierre}
\begin{tabular}{@{}>{\raggedright\arraybackslash}p{0.22\linewidth}>{\raggedright\arraybackslash}p{0.34\linewidth}>{\raggedright\arraybackslash}p{0.34\linewidth}@{}}
\toprule
Aspect & Alcubierre~\cite{alcubierre1994} & PGA--DST (this work) \\
\midrule
Mechanism &
Contracting/expanding spacetime bubble via $g_{\mu\nu}$ &
Collective dual-rotor phase tuning; $J<0$ layers \\
Exotic matter &
Required; violates NEC~\cite{ford2000} &
Not required; $J$ sign from Killing-form asymmetry \\
Causality &
Local $c$ exceeded inside bubble &
Effective superluminal \emph{coast}; local $c$ preserved via $M\widetilde{M}=1$ \\
Information &
Can outrun light within bubble &
Bound to light cone; $\eta_{\mu\nu}\lambda^\mu\lambda^\nu=0$ \\
GR equivalence &
Modified Einstein equations &
TEGR/EH recovered; same weak-field limits \\
Engineering &
Negative energy density at bubble wall &
THz circularly polarised driving of $\theta_a$ \\
\bottomrule
\end{tabular}
\end{table}

The contrast is ontological as well as operational. Alcubierre engineering deforms a background metric. PGA--DST has no background continuum; it modulates the torsion scalar $J$ and, during coast, the translational component of $J^{(5)}$ on approximately light-like worldlines.

\section{Repulsive Layers}
\label{sec:repulsive}

\subsection{Sign structure of J}

\begin{table}[ht]
\centering
\caption{Torsion scalar sign and physical interpretation}
\label{tab:J-signs}
\begin{tabular}{@{}cll@{}}
\toprule
Condition & Sign of $J$ & Effect \\
\midrule
$\alpha_a \gg \beta_a$ (aligned) & $J > 0$ & Attractive gravity \\
$\alpha_a \approx \beta_a$ & $J \approx 0$ & Vacuum / shielding / inertia reduction \\
$\beta_a > \alpha_a$ (repulsive point) & $J < 0$ & Repulsive gravity \\
\bottomrule
\end{tabular}
\end{table}

Repulsive layers occur near $\theta_a=(2n+1)\pi$. In strong-field astrophysical contexts, the same sign oscillation produces stratified ``onion'' structures (twist stars) rather than event horizons and singularities~\cite{dual2025}: an infinite stack of alternating $J>0$ and $J<0$ shells provides the cosmological backdrop for layer-hopping.

In any region $\mathcal{R}$ with $J(x)<0$, dual-sector rotation dominates, yielding effective repulsion. Such repulsive layers are engineerable in principle via coherent excitation of $\theta_a$. For warp entry, a $J<0$ shell is not merely a barrier but the \emph{operational medium}: the vessel transitions into a locally generated or cosmically encountered repulsive layer, coasts with reduced gravitational coupling, and exits at a scheduled phase boundary. Stability under perturbation follows from self-reinforcing mismatch dynamics: intrusion of external matter increases $|\beta_a|$, driving $J$ more negative.

\section{Drive Mechanism: THz Helicity and Dual-Rotor Coupling}
\label{sec:drive}

\subsection{Electromagnetic splitting}

The electromagnetic field splits across DST sectors:
\begin{equation}
F = F_{\mathrm{usual}} + F_{\mathrm{dual}},
\qquad
F_{\mathrm{usual}} = \mathbf{E}\cdot(iI,iJ,iK),
\qquad
F_{\mathrm{dual}} = \mathbf{B}\cdot(I,J,K).
\label{eq:F-split}
\end{equation}
Electric fields couple to hyperbolic (usual) generators; magnetic fields couple exclusively to cyclic (dual) generators. This algebraic separation is the geometric origin of the distinct transformation properties under time reversal ($\mathbf{E}\to\mathbf{E}$, $\mathbf{B}\to-\mathbf{B}$).

\subsection{Circular polarisation as a chiral switch}

For a circularly polarised wave with helicity $\sigma=\pm 1$ propagating along $\hat{z}$,
\begin{equation}
\langle F_{\mathrm{dual}}\rangle \propto \sigma\,E_0^2\,K.
\label{eq:Fdual-avg}
\end{equation}
Helicity directly controls $J$:
\begin{itemize}[leftmargin=2em]
  \item $\sigma=+1$ (right-handed): enhances attractive torsion ($J>0$).
  \item $\sigma=-1$ (left-handed): reduces $J$; sufficient intensity can invert to $J<0$.
\end{itemize}
The resonance driving equation is
\begin{equation}
\dot{\theta} = \kappa\,\sigma\,E_0^2\sin(\omega t - \theta + \delta_\sigma),
\label{eq:theta-dot}
\end{equation}
where $\kappa$ is a matter-dependent coupling constant. Synchronising $\theta$ with $\phi$ drives $J\to 0$ (shielding, inertia reduction); driving $\theta$ past $\phi$ yields $J<0$ (repulsion, warp entry).

\subsection{THz collective resonance}

Single-particle dual resonance lies at Planckian frequencies ($\sim 10^{43}\,\mathrm{Hz}$). Macroscopic coherence over $N\sim 10^{26}$--$10^{32}$ particles redshifts the effective resonance to the THz band:
\begin{equation}
\omega_{\mathrm{eff}} \sim \frac{\omega_{\mathrm{Planck}}}{\sqrt{N}} \sim 10^{12}\text{--}10^{15}\,\mathrm{Hz}.
\label{eq:freq-redshift}
\end{equation}
Superconducting phase locking and high-$Q$ cavities further enhance coupling. This is the same mechanism proposed for table-top tests of $\Delta m/m\gtrsim 10^{-6}$ with helicity-reversal signatures~\cite{dual2025}. Warp engineering is the macroscopic extrapolation of that laboratory programme.

\section{Operational Regimes: On-Demand Entry vs Layer-Hopping}
\label{sec:operational}

\subsection{Cosmic phase drift}

Large-scale structure induces slow background drift of dual-rotor phases:
\begin{equation}
\left|\frac{d\theta_a}{dl}\right| \simeq 1.1\times 10^{-26}~\mathrm{rad\cdot m^{-1}},
\qquad
\frac{d\theta_a}{dt} \simeq 3.3\times 10^{-10}\,(v/c)~\mathrm{rad\cdot s^{-1}}.
\label{eq:drift}
\end{equation}
This continuously sweeps the infinite stack of repulsive layers past an observer moving at velocity $v$.

\subsection{Effective Hamiltonian and entry energy}

Near a repulsive layer, with $\theta=\theta_a-(2n+1)\pi$,
\begin{equation}
H = \frac{p^2}{2m_{\mathrm{eff}}} + V(\theta),
\qquad
V(\theta) \approx \tfrac{1}{2}k\theta^2 - \tfrac{\hbar\omega_0}{2}\cos(\theta),
\label{eq:hamiltonian}
\end{equation}
where $m_{\mathrm{eff}}$ is the collective dual excitation inertia and $\omega_0\sim 10^{14}$--$10^{15}\,\mathrm{Hz}$ for THz collective modes. The semiclassical energy to reach phase offset $\Delta\phi$ is
\begin{equation}
E(\Delta\phi) \approx N_b\,\hbar\omega_{\mathrm{res}}\,\Bigl|\sin\!\Bigl(\tfrac{\Delta\phi}{2}\Bigr)\Bigr|,
\label{eq:energy}
\end{equation}
with $N_b\approx 6\times 10^{32}$ baryons in a $10^6\,\mathrm{kg}$ vessel. Small $\Delta\phi$ gives quadratic scaling; near $\pi$ the cost saturates, explaining the extreme sensitivity of on-demand entry.

During entry, the motor \eqref{eq:motor} is driven primarily through $\Omega_{\mathrm{torsion}}$ toward $J<0$. During coast in a repulsive layer, translational mismatch $\lambda^\mu$ is tuned to lie near the light cone so that the $J^{(5)}$ translational term is suppressed and displacement proceeds via phase-layer sliding rather than acausal local boosts.

\subsection{On-demand entry energy}

Worst-case on-demand entry ($\Delta\phi_{\max}\approx\pi$) for a $10^6\,\mathrm{kg}$ vessel:

\begin{table}[ht]
\centering
\caption{On-demand subspace entry (worst-case $\Delta\phi\simeq\pi$)}
\label{tab:ondemand}
\begin{tabular}{@{}lcr@{}}
\toprule
Cruise velocity & $\gamma$ & Energy required ($10^6$ kg vessel) \\
\midrule
$0.10c$ & 1.005 & $\sim 8\times 10^{11}\,\mathrm{J}\approx\qty{220}{GWh}$ \\
$0.30c$ & 1.048 & $\sim 7.5\times 10^{11}\,\mathrm{J}\approx\qty{210}{GWh}$ \\
$0.50c$ & 1.155 & $\sim 6.9\times 10^{11}\,\mathrm{J}\approx\qty{190}{GWh}$ \\
$0.99c$ & $\sim 7$ & $\sim 3\times 10^{11}\,\mathrm{J}\approx\qty{80}{GWh}$ \\
\bottomrule
\end{tabular}
\end{table}

These values approach projected compact fusion or antimatter-catalysed limits, making pure on-demand entry technologically challenging.

\subsection{Layer-hopping performance}

Waiting time for phase deficit below threshold $\Delta\phi_{\mathrm{threshold}}$:
\begin{equation}
\tau \simeq \frac{\Delta\phi_{\mathrm{threshold}}}{|d\theta/dt|}
= 3\times 10^9\,\Delta\phi_{\mathrm{threshold}}\,\Bigl(\frac{c}{v}\Bigr)~\mathrm{s}.
\label{eq:wait}
\end{equation}

\begin{table}[ht]
\centering
\caption{Layer-hopping performance at $v=0.3c$ for a $10^6$ kg vessel}
\label{tab:layerhop}
\begin{tabular}{@{}lcccc@{}}
\toprule
Threshold $\Delta\phi$ & Typical wait & Energy required & Power (10-min pulse) & Gain \\
\midrule
$\pi$ (on-demand) & 0 & \qty{750}{GWh} & --- & $1\times$ \\
$10^{-1}$ & $\sim$1 yr & \qty{75}{GWh} & \qty{450}{MW} & $\sim 10^4\times$ \\
$10^{-2}$ & $\sim$40 days & \qty{7.5}{GWh} & \qty{45}{MW} & $\sim 10^5\times$ \\
$10^{-3}$ & $\sim$4 days & \qty{750}{MWh} & \qty{4.5}{MW} & $\sim 10^6\times$ \\
$10^{-4}$ & $\sim$10 hrs & \qty{75}{MWh} & \qty{450}{kW} & $\sim 10^7\times$ \\
$10^{-5}$ & $\sim$1 hr & \qty{7.5}{MWh} & \qty{45}{kW} & $\sim 10^8\times$ \\
\bottomrule
\end{tabular}
\end{table}

A threshold $\Delta\phi\approx 10^{-3}$ offers a practical compromise between waiting time and power. For a 10 light-year mission, total energy expenditure drops to 10--100 MWh under the hybrid protocol.

\section{Causality, Horizons, and Effective Superluminal Coasting}
\label{sec:causality}

\subsection{Unitary motors and local causality}

Theorem~\ref{thm:unitarity} guarantees that every physical transformation is implemented by a unitary motor $M\widetilde{M}=1$. No sandwich product can produce acausal local propagation: the light cone $\eta_{\mu\nu}\lambda^\mu\lambda^\nu=0$ is built into $J^{(5)}$ \eqref{eq:J5}. Photons correspond to full anti-synchronisation ($\phi=\theta$, $\Omega=1$) with vanishing translational contribution in $J^{(5)}$.

\subsection{Effective superluminal displacement}

``Effective superluminal coasting'' means that the \emph{phase-layer displacement rate} of the vessel exceeds $c$ as measured against the ambient $J>0$ reference frame, while:
\begin{enumerate}[leftmargin=2em]
  \item local signal velocity remains $\le c$ in the vessel's instantaneous frame;
  \item information about the hop is carried only on the light channel;
  \item no closed timelike curves are generated because hops are finite-duration transitions between globally hyperbolic TEGR patches.
\end{enumerate}

This is analogous to a refractive index change in a medium: phase velocity can exceed $c$ without violating relativity, provided group velocity and information velocity do not.

\subsection{Horizons and quasi-horizons}

PGA--DST excludes true event horizons and central singularities in favour of stratified twist-star interiors~\cite{dual2025}. However, \emph{quasi-horizons}---regions where $\int J\,dr$ rises steeply---can impede exit without supplying sufficient torsion energy. Warp navigation must map local phase gradients to avoid trapping in quasi-horizon shells. This is consistent with the layer-stability analysis below.

\section{Mission Protocol, Stability, and Navigation}
\label{sec:mission}

\subsection{Hybrid interstellar protocol}

\begin{enumerate}[leftmargin=2em]
  \item Cruise at $0.2$--$0.5c$ using nuclear or beamed propulsion.
  \item Monitor local dual phase via a compact THz resonant cavity coupled to dual degrees of freedom.
  \item When $\Delta\phi<10^{-3}\,\mathrm{rad}$ (every few days to weeks at $0.3c$), apply a low-power, left-handed circularly polarised THz pulse to complete the transition into $J<0$.
  \item Coast in the repulsive layer for weeks to decades, achieving effective superluminal displacement.
  \item Exit via complementary pulse (helicity reversal or phase realignment) and repeat, scheduling hops from phase-gradient maps.
\end{enumerate}

\subsection{Layer stability}

Inside a repulsive layer, small deviations $\delta\theta$ experience restoring forces from $V(\theta)$ in \eqref{eq:hamiltonian}. Characteristic oscillation frequency:
\begin{equation}
\omega_{\mathrm{stab}} \approx \sqrt{\frac{k}{m_{\mathrm{eff}}}}.
\end{equation}
Active feedback via the same THz cavity damps unwanted excitations. Self-reinforcing mismatch dynamics in $J<0$ regions provide additional stability against external perturbation.

\subsection{Navigation in torsional phase space}

Navigation requires control of entry/exit points and hop direction. We propose auxiliary low-mass probes to map local $\theta_a$ gradients, enabling predictive scheduling analogous to gravitational slingshots but in torsional phase space. Motor factorisation $M=RT$ separates rotational mismatch (layer selection via $R$) from translational coast (via light-like $T$).

\section{Experimental Roadmap and Open Questions}
\label{sec:roadmap}

\subsection{Near-term laboratory tests}

Before any warp application, the foundational hypothesis---macroscopic dual-phase coherence coupled to EM fields---must be tested:
\begin{enumerate}[leftmargin=2em]
  \item \textbf{Table-top gravity engineering.} Superconducting test masses in high-$Q$ THz cavities; monitor $\Delta m/m\gtrsim 10^{-6}$ with helicity-reversal sign change~\cite{dual2025}.
  \item \textbf{Repulsive-layer generation.} High-intensity circularly polarised THz ($10^{12}$--$10^{15}\,\mathrm{Hz}$) to drive $J<0$ locally; measure repulsive deflection or weight reduction.
  \item \textbf{Resonance spectroscopy.} Systematic frequency/helicity scans to constrain $\kappa$ and compactification scale $r_{\mathrm{comp}}$.
\end{enumerate}

\subsection{Astronomical and space-based tests}

\begin{enumerate}[leftmargin=2em]
  \item Precision torsion balances or atom interferometers to detect cosmic phase drift \eqref{eq:drift}.
  \item Satellite baselines to measure large-scale dual-phase gradients across the solar system.
  \item Correlation of repulsive-layer signatures with galactic-scale $J<0$ boundaries in rotation-curve data~\cite{dual2025}.
\end{enumerate}

\subsection{Key open questions}

\begin{itemize}[leftmargin=2em]
  \item What is the precise coupling $\kappa$ between EM fields and dual rotors?
  \item What is the decoherence time for macroscopic dual-phase coherence?
  \item Can collective $N\sim 10^{32}$ baryon locking be maintained during a hop pulse?
  \item How do quasi-horizon shells interact with repeated layer transitions?
\end{itemize}

\section{Discussion and Conclusions}
\label{sec:conclusions}

We have upgraded the dual spacetime warp propulsion hypothesis to the PGA--DST framework. The central claims are:

\begin{enumerate}[leftmargin=2em]
  \item Warp is resonance tuning of collective dual rotors, not Alcubierre metric engineering.
  \item Repulsive layers ($J<0$) provide the operational medium for effective superluminal coasting.
  \item THz circularly polarised driving with helicity $\sigma=-1$ is the proposed entry mechanism.
  \item Layer-hopping exploits cosmic phase drift to reduce entry energy by 4--7 orders of magnitude relative to on-demand entry.
  \item Motor unitarity $M\widetilde{M}=1$ and the light-cone structure of $J^{(5)}$ preserve local causality.
\end{enumerate}

All quantitative estimates inherit the semiclassical model \eqref{eq:energy} and remain \emph{hypothetical} until macroscopic dual-phase coherence is experimentally confirmed. Failure of table-top tests would constrain $\kappa$ and $r_{\mathrm{comp}}$ without falsifying the algebraic PGA--DST structure, which independently recovers GR via TEGR. Success at laboratory scales would open a path toward engineering repulsive layers and, in principle, interstellar layer-hopping.

\begin{thebibliography}{99}

\bibitem{alcubierre1994}
M.~Alcubierre,
``The warp drive: hyper-fast travel within general relativity,''
\emph{Classical and Quantum Gravity} \textbf{11}, L73 (1994).

\bibitem{ford2000}
L.~H.~Ford and T.~A.~Roman,
``Negative energy, wormholes and warp drive,''
\emph{Scientific American} \textbf{281}, 46 (1999);
see also arXiv:gr-qc/0001073.

\bibitem{dual2025}
Anonymous,
``Gravity as Torsion between Dual Spacetime: A Biquaternionic Reformulation of General Relativity,''
December 2025 (in preparation).

\bibitem{pga2025}
Anonymous,
``\(G(3,1,1)\) Extension of Double Spacetime Torsional Mismatch: 10-Dimensional Generator Structure and Plane-Based Bracket Products,''
2025.

\bibitem{tegr}
S.~Bahamonde \emph{et al.},
``Teleparallel Gravity: A Review,''
arXiv:2106.13757 (2021).

\bibitem{Gunn2017}
C.~Gunn,
``Geometric Algebra for Computer Graphics,''
Springer (2017).

\bibitem{Lengyel2024}
V.~Lengyel,
``Projective Geometric Algebra,''
BCS (2024).

\bibitem{Selig2005}
J.~M.~Selig,
``Geometric Fundamentals of Robotics,''
2nd ed., Springer (2005).

\bibitem{DoranLasenby2003}
C.~Doran and A.~Lasenby,
``Geometric Algebra for Physicists,''
Cambridge University Press (2003).

\end{thebibliography}

\end{document}
